% GENERATED FILE. Edit canonical lessons, then run npm run generate:books.
% canonical-source-hash: fnv1a64:889021445088f95c
% canonical-lessons: GU-C220-vahevu, GU-C220-chalkavu, GU-C220-pur, GU-C220-osarvu, GU-C220-ughadvu

\chapter{When the river rises}
\label{ch:gu-when-the-river-rises}
\hlchaptermodality{220}

\begin{chapteropening}
\textbf{By the end of this chapter:} \emph{I can say that the river will overflow and that the water will go down.}

Complete the last lesson of chapter 220: I can say that the river will overflow and that the water will go down.
\end{chapteropening}

\section[\texorpdfstring{\gu{વહેવું} (vahevũ)}{વહેવું (vahevũ)}]{\gu{વહેવું} (vahevũ) — to flow}

\label{lesson:GU-C220-vahevu}



\begin{quote}
Before the new one: say the Gujarati for hailstones, then the Gujarati for unseasonal rain.
\end{quote}

\subsection*{You'll want to know: \gu{વહેવું}}
\textbf{\gu{વહેવું}} — \emph{vahevũ} — \textquotedblleft{}to flow\textquotedblright{}.

\begin{quote}
\textbf{\gu{વરસાદ} \gu{પછી} \gu{નદી} \gu{વહેશે}.} — \emph{varsād pachī nadī vaheśe} — \textquotedblleft{}After the rain the river will flow.\textquotedblright{}
\end{quote}

\begin{cousinweb}
\textbf{\gu{વહેવું}} (\emph{vahevũ}) is a tadbhava. It comes from Sanskrit \emph{vahati}, \textquotedblleft{}carries, bears along (of rivers)\textquotedblright{}, and CDIAL glosses the Gujarati verb \textquotedblleft{}to flow\textquotedblright{} (CDIAL 11453).
\end{cousinweb}

\begin{grammarlens}[title={What you've built}]
One more word for everyday life at A2.
\end{grammarlens}

\subsection*{Your turn}
\begin{itemize}
\raggedright
  \item \emph{Say it:} \emph{vahevũ}
  \item \emph{Say it:} \emph{vahevũ}, once more
  \item \emph{Say it:} \emph{varsād pachī nadī vaheśe} — After the rain the river will flow.
  \item \emph{Recall:} say the Gujarati for hailstones, then the Gujarati for unseasonal rain, then say \emph{vahevũ} again
\end{itemize}

\begin{tcolorbox}[breakable,colback=teal!4,colframe=teal!35!black,title={Before you move on}]
What does \gu{વહેવું} mean? (\textquotedblleft{}To flow\textquotedblright{}.) Say it once more.
\end{tcolorbox}
\section[\texorpdfstring{\gu{છલકાવું} (chalkāvũ)}{છલકાવું (chalkāvũ)}]{\gu{છલકાવું} (chalkāvũ) — to overflow, to brim over}

\label{lesson:GU-C220-chalkavu}



\begin{quote}
Before the new one: say the Gujarati for unseasonal rain, then the Gujarati for to flow.
\end{quote}

\subsection*{You'll want to know: \gu{છલકાવું}}
\textbf{\gu{છલકાવું}} — \emph{chalkāvũ} — \textquotedblleft{}to overflow, to brim over\textquotedblright{}.

\begin{quote}
\textbf{\gu{કાલે} \gu{તળાવ} \gu{છલકાશે}.} — \emph{kāle taḷāv chalkāśe} — \textquotedblleft{}The lake will overflow tomorrow.\textquotedblright{}
\end{quote}

\begin{cousinweb}
\textbf{\gu{છલકાવું}} (\emph{chalkāvũ}) comes from \emph{*chala}, \textquotedblleft{}splash\textquotedblright{}, an imitative reconstructed form; CDIAL lists the Gujarati verb, glossed \textquotedblleft{}to be spilt\textquotedblright{} (CDIAL 5002).
\end{cousinweb}

\begin{grammarlens}[title={What you've built}]
One more word for everyday life at A2.
\end{grammarlens}

\subsection*{Your turn}
\begin{itemize}
\raggedright
  \item \emph{Say it:} \emph{chalkāvũ}
  \item \emph{Say it:} \emph{chalkāvũ}, once more
  \item \emph{Say it:} \emph{kāle taḷāv chalkāśe} — The lake will overflow tomorrow.
  \item \emph{Recall:} say the Gujarati for unseasonal rain, then the Gujarati for to flow, then say \emph{chalkāvũ} again
\end{itemize}

\begin{tcolorbox}[breakable,colback=teal!4,colframe=teal!35!black,title={Before you move on}]
What does \gu{છલકાવું} mean? (\textquotedblleft{}To overflow, to brim over\textquotedblright{}.) Say it once more.
\end{tcolorbox}
\section[\texorpdfstring{\gu{પૂર} (pūr)}{પૂર (pūr)}]{\gu{પૂર} (pūr) — flood}

\label{lesson:GU-C220-pur}



\begin{quote}
Before the new one: say the Gujarati for to flow, then the Gujarati for to overflow.
\end{quote}

\subsection*{You'll want to know: \gu{પૂર}}
\textbf{\gu{પૂર}} — \emph{pūr} — \textquotedblleft{}flood\textquotedblright{}.

\begin{quote}
\textbf{\gu{નદીમાં} \gu{પૂર} \gu{આવશે}.} — \emph{nadīmā̃ pūr āvśe} — \textquotedblleft{}The river will flood.\textquotedblright{}
\end{quote}

A neuter noun.

\begin{cousinweb}
\textbf{\gu{પૂર}} (\emph{pūr}) is inherited from Sanskrit \emph{pūra}, \textquotedblleft{}filling\textquotedblright{}, which already meant \textquotedblleft{}flood\textquotedblright{} in the medical texts (CDIAL 8330).
\end{cousinweb}

\begin{grammarlens}[title={What you've built}]
One more word for everyday life at A2.
\end{grammarlens}

\subsection*{Your turn}
\begin{itemize}
\raggedright
  \item \emph{Say it:} \emph{pūr}
  \item \emph{Say it:} \emph{pūr}, once more
  \item \emph{Say it:} \emph{nadīmā̃ pūr āvśe} — The river will flood.
  \item \emph{Recall:} say the Gujarati for to flow, then the Gujarati for to overflow, then say \emph{pūr} again
\end{itemize}

\begin{tcolorbox}[breakable,colback=teal!4,colframe=teal!35!black,title={Before you move on}]
What does \gu{પૂર} mean? (\textquotedblleft{}Flood\textquotedblright{}.) Say it once more.
\end{tcolorbox}
\section[\texorpdfstring{\gu{ઓસરવું} (osarvũ)}{ઓસરવું (osarvũ)}]{\gu{ઓસરવું} (osarvũ) — to recede, to go down}

\label{lesson:GU-C220-osarvu}



\begin{quote}
Before the new one: say the Gujarati for to overflow, then the Gujarati for flood.
\end{quote}

\subsection*{You'll want to know: \gu{ઓસરવું}}
\textbf{\gu{ઓસરવું}} — \emph{osarvũ} — \textquotedblleft{}to recede, to go down\textquotedblright{}.

\begin{quote}
\textbf{\gu{બે} \gu{દિવસમાં} \gu{પાણી} \gu{ઓસરશે}.} — \emph{be divasmā̃ pāṇī osarśe} — \textquotedblleft{}The water will go down in two days.\textquotedblright{}
\end{quote}

\begin{cousinweb}
\textbf{\gu{ઓસરવું}} (\emph{osarvũ}) is a tadbhava. It comes from Sanskrit \emph{apasarati}, \textquotedblleft{}slips off from, goes away\textquotedblright{}, and CDIAL glosses the Gujarati verb \textquotedblleft{}to recede, decrease, subside\textquotedblright{} (CDIAL 465).
\end{cousinweb}

\begin{grammarlens}[title={What you've built}]
One more word for everyday life at A2.
\end{grammarlens}

\subsection*{Your turn}
\begin{itemize}
\raggedright
  \item \emph{Say it:} \emph{osarvũ}
  \item \emph{Say it:} \emph{osarvũ}, once more
  \item \emph{Say it:} \emph{be divasmā̃ pāṇī osarśe} — The water will go down in two days.
  \item \emph{Recall:} say the Gujarati for to overflow, then the Gujarati for flood, then say \emph{osarvũ} again
\end{itemize}

\begin{tcolorbox}[breakable,colback=teal!4,colframe=teal!35!black,title={Before you move on}]
What does \gu{ઓસરવું} mean? (\textquotedblleft{}To recede, to go down\textquotedblright{}.) Say it once more.
\end{tcolorbox}
\section[\texorpdfstring{\gu{ઊઘડવું} (ūghaḍvũ)}{ઊઘડવું (ūghaḍvũ)}]{\gu{ઊઘડવું} (ūghaḍvũ) — to clear up, to open}

\label{lesson:GU-C220-ughadvu}



\begin{quote}
Before the new one: say the Gujarati for flood, then the Gujarati for to recede.
\end{quote}

\subsection*{You'll want to know: \gu{ઊઘડવું}}
\textbf{\gu{ઊઘડવું}} — \emph{ūghaḍvũ} — \textquotedblleft{}to clear up, to open\textquotedblright{}.

\begin{quote}
\textbf{\gu{વરસાદ} \gu{પછી} \gu{આકાશ} \gu{ઊઘડશે}.} — \emph{varsād pachī ākāś ūghaḍśe} — \textquotedblleft{}The sky will clear after the rain.\textquotedblright{}
\end{quote}

\begin{cousinweb}
\textbf{\gu{ઊઘડવું}} (\emph{ūghaḍvũ}) is a tadbhava. It comes from \emph{*udghaṭati}, \textquotedblleft{}is open\textquotedblright{}, and CDIAL glosses the Gujarati verb \textquotedblleft{}to open, bloom, be free from rain\textquotedblright{} (CDIAL 1968).
\end{cousinweb}

\begin{grammarlens}[title={What you've built}]
One more word for everyday life at A2.
\end{grammarlens}

\subsection*{Your turn}
\begin{itemize}
\raggedright
  \item \emph{Say it:} \emph{ūghaḍvũ}
  \item \emph{Say it:} \emph{ūghaḍvũ}, once more
  \item \emph{Say it:} \emph{varsād pachī ākāś ūghaḍśe} — The sky will clear after the rain.
  \item \emph{Recall:} say the Gujarati for flood, then the Gujarati for to recede, then say \emph{ūghaḍvũ} again
\end{itemize}

\begin{tcolorbox}[breakable,colback=teal!4,colframe=teal!35!black,title={Before you move on}]
What does \gu{ઊઘડવું} mean? (\textquotedblleft{}To clear up, to open\textquotedblright{}.) Say it once more.
\end{tcolorbox}
